\documentclass[11pt]{article}
\usepackage[T1]{fontenc}
\usepackage[utf8]{inputenc}
\usepackage{lmodern,amsmath,amssymb,amsthm,mathtools}
\usepackage[a4paper,margin=27mm]{geometry}
\usepackage{booktabs,longtable,array,enumitem,microtype,xurl}
\usepackage[colorlinks=true,linkcolor=blue,citecolor=blue,urlcolor=blue]{hyperref}
\usepackage{fancyhdr}
\pagestyle{fancy}\fancyhf{}\fancyhead[L]{THETA3: three signatures}\fancyhead[R]{21 November 2026}\fancyfoot[C]{\thepage}
\setlength{\headheight}{14pt}
\newcommand{\Q}{\mathbb Q}\newcommand{\C}{\mathbb C}\newcommand{\Z}{\mathbb Z}\newcommand{\R}{\mathbb R}\newcommand{\A}{\mathbb A}\newcommand{\OO}{\mathcal O}
\newcommand{\Hdg}{\operatorname{Hdg}}\newcommand{\CH}{\operatorname{CH}}\newcommand{\Hom}{\operatorname{Hom}}\newcommand{\Tr}{\operatorname{Tr}}\newcommand{\SU}{\operatorname{SU}}\newcommand{\U}{\operatorname{U}}
\newcommand{\im}{\operatorname{im}}\newcommand{\vol}{\operatorname{vol}}\newcommand{\Res}{\operatorname{Res}}\newcommand{\Nm}{\operatorname{N}}\newcommand{\cl}{\operatorname{cl}}\newcommand{\Span}{\operatorname{span}}
\newcommand{\framework}[1]{\par\smallskip\noindent\textit{Framework and scope.} #1\par}
\newtheorem{theorem}{Theorem}[section]\newtheorem{proposition}[theorem]{Proposition}\newtheorem{lemma}[theorem]{Lemma}\newtheorem{corollary}[theorem]{Corollary}
\theoremstyle{definition}\newtheorem{definition}[theorem]{Definition}\newtheorem{remark}[theorem]{Remark}
\title{THETA3\\Theta kernels, definite periods and integral models\\\large Exact comparisons and quantitative certificates}
\author{Run THETA3-20261121\\GPT-6 Astra Pro, with independent agent proof checks}
\date{21 November 2026}
\begin{document}\maketitle
\begin{abstract}
We compare the supplied CM Hodge construction with the theta and period mechanisms in Alp\"oge's archive. The comparison keeps the Hermitian spaces, splitting characters, coefficient types, Haar measures and integral structures explicit. A zero-orbit identity yields a quantitative finite search for nonzero mixed theta periods. A totally definite rank-two subgroup inside the original ball group yields a restriction seesaw and finite-dimensional recovery of the ball pairing, with a positive correction operator. Local calculations identify a canonical Bessel value, the limitations of signature replacement, and the extra information needed to transfer determinant bounds. A noncongruence cover produces a non-CM/CM mixed period, while a projector-and-rank condition isolates the remaining Hodge extraction problem. In the integral setting we calculate both negative derived intersections and a positive rank-one sector, and derive the height-two formal derivative slope occurring in Alp\"oge's step~2. Every new statement is separated from cited background and from unresolved identifications. No claim depending on Masullo is used.
\end{abstract}
\noindent\textbf{Working status.} This is the evolving research paper for the requested run. Checked mathematical statements below have the scope stated in their hypotheses; an unchecked global identification is never used as a theorem. The final status and current gates are recorded in \texttt{QUESTIONS.tsv} and \texttt{LEDGER.md}.

\tableofcontents\clearpage
\part{Scope, sources and the geometric dictionary}
\section{What is being compared}
The requested sketch is a map of proposed connections, not a collection of background theorems. In particular, sketch~1.7 does not specify the test vectors, splittings, inner forms or normalizing constants of its ``one seesaw pairing.'' We therefore distinguish exact identities, exact nonvanishing implications, and a still-uninstantiated identification with the character-independent iterated functional in Alp\"oge's step~12. This is a mathematical distinction: a local Bessel integral, a global period on a definite Shimura set, and a cup-product integral on a surface have different domains.

Throughout, $E/F$ is a CM extension with $F$ totally real. A skew-Hermitian multiplicity space is denoted $W$; writing $W=\delta P$ for a Hermitian space $P$ requires a fixed $\delta\in E^\times$ with $\bar\delta=-\delta$. The associated symplectic form is the trace of the tensor-product form. A change of $\delta$, an additive character, a Weil splitting or a determinant-line trivialization can change a numerical period. These choices are not suppressed.

In the surface construction, $V$ is Hermitian of dimension~3, with signature $(2,1)$ at one real place $\mu$ and $(3,0)$ at the others. There is at least one definite place. Thus $V$ is anisotropic over $F$: an isotropic rational vector would remain isotropic at a definite real place. A sufficiently small level gives a compact projective ball surface. We distinguish $\SU(V)$, used in the original geometric construction, from $\U(V)$, used for adelic theta identities. Passing between their appropriate components and finite covers retains determinant characters and component measures.

\subsection{Background accepted and what was checked}
The supplied 032 theorem for CM abelian varieties, the supplied human proof of van Hoften--Sempliner, and the stated Alp\"oge RH and Dirichlet GRH theorems are carried as cited background, as requested. This run does not purport to referee those global theorems. It reads and checks the particular hypotheses and formulas used. No Masullo claim enters any deduction.

The supplied ZIP contains two Lean entry files and a source map, not the full tree. Named modules were subsequently obtained from a public mirror at commit
\texttt{7cbf7a0c32fca1badf05931879a37c3f8a2b03b4}.
The mirror's \texttt{Main.lean} and \texttt{Statements.lean} match the supplied files byte for byte. This does not authenticate every mirror module against the original multi-part archive. No Lean compilation was performed. A Lean \texttt{def ... : Prop} is a definition of a proposition; it is not itself a proof of that proposition. Source readings and elementary consequences of those definitions are recorded accordingly.

A mathematical finding here means a deduction formulated and proved for this comparison; it does not mean a claim of research priority. Standard consequences included to prevent an incorrect inference are marked as such. The independent checks are whole-proof readings or reconstructions, not independent formal certifications. The model was GPT-6 Astra Pro; no numerical reasoning-effort setting was exposed. Actual check events, including background uses, appear in the TSV register.

\subsection{Realizations and correspondences}
A nonzero mixed period is a nonzero functional on a Betti realization. It produces an algebraic cycle only after an algebraic morphism or correspondence has been constructed and an extraction criterion has been verified. In the 032 setting rational maps from CM Hodge structures into $H^1(X,\Q)$ come from maps to abelian varieties after multiplying by an integer. Pushing forward $[X]$ then gives the surface cycle. An analytic inverse, an evaluation of a differential at a CM point, or a positive height is not by itself such an algebraic correspondence.

We use motives modulo homological equivalence when discussing projectors whose idempotence is proved on cohomology. We do not silently assert idempotence in the Chow ring. A homologically trivial cycle is numerically trivial, although it can carry nonzero Abel--Jacobi or height data. These distinctions are essential in comparing the ordinary and arithmetic theta correspondences.

\subsection{Grothendieck anchors}
The archival anchors in this paper locate the framework of a finding, rather than claim that its theta formula appears in Grothendieck's text. The original scans of cote~44, scans~7--8, were read: the letter of 12 August 1964 discusses the intermediate Jacobians $J^i$, their duality and the correspondence pairing used to reduce positivity to a polarization question. Cote~16, scans~43--46, the letter to Coates dated 4 January 1967, discusses the standard conjectures, equality of dimensions of complementary algebraic classes, and algebraic inverses associated with hyperplane sections. In particular the date and the first dimension statement are on scan~43, not scan~44. The published \emph{Standard conjectures on algebraic cycles}, pp.~196--198, supplies the complementary realization/polarization language. Later sections give verified SGA/EGA locators for specialization, cycle classes and spreading out.

No separate ``text of record'' transcription of these archival pages was supplied. Consequently a typographical problem in extracted modern-paper text is not entered as a correction to Grothendieck's page. The transcription register records only actual comparisons for which both readings exist.
